% Generated draft: requires mathematical and source review.
\documentclass[11pt]{article}
\usepackage[margin=25mm]{geometry}
\usepackage{fontspec}
\setmainfont{Latin Modern Roman}
\usepackage{amsmath,amssymb,mathtools,unicode-math}
\setmathfont{Latin Modern Math}
\usepackage{enumitem,needspace,xurl,hyperref}
\setlength{\parindent}{0pt}
\setlength{\parskip}{5pt}
\setlength{\emergencystretch}{3em}
\newcommand{\N}{\mathbb{N}}
\newcommand{\Z}{\mathbb{Z}}
\newcommand{\Q}{\mathbb{Q}}
\newcommand{\R}{\mathbb{R}}
\newcommand{\C}{\mathbb{C}}
\newcommand{\F}{\mathbb{F}}
\newcommand{\A}{\mathbb{A}}
\newcommand{\PP}{\mathbb{P}}
\newcommand{\E}{\mathbb{E}}
\newcommand{\eps}{\varepsilon}
\newcommand{\dd}{\,\mathrm{d}}
\begin{document}
\section*{UnsolvedMath ID: 30001332}
\textbf{Source problem code:} catalogue entry
\subsection{Definitions and mathematical statement}
The supplied question is: ``Are homotopic projections in matrices over a complex group ring similar?'' Its precise interpretation requires conventions absent from the supplied record.

For a group \(G\), its complex group ring is
\[
\C[G]=\left\{\sum_{g\in G}a_g g : a_g\in\C,\ a_g=0\text{ for all but finitely many }g\right\},
\]
with multiplication induced by the group law. Its customary involution is
\[
\left(\sum_{g\in G}a_g g\right)^*=\sum_{g\in G}\overline{a_g}\,g^{-1}.
\]
For a positive integer \(n\), this induces the conjugate-transpose involution on \(M_n(\C[G])\). Under the self-adjoint convention, a projection is a matrix \(p\) satisfying \(p^2=p=p^*\). If the source instead uses ``projection'' for an arbitrary idempotent, only \(p^2=p\) is required; the supplied record does not settle this convention.

For \(p,q\in M_n(\C[G])\), similarity over the group ring means
\[
\exists u\in\operatorname{GL}_n(\C[G])\qquad q=upu^{-1},
\]
where \(\operatorname{GL}_n(\C[G])\) denotes the matrices invertible over \(\C[G]\).

To define homotopy by paths, one must specify a topology \(\tau\) on the matrix algebra. Under a same-size, internal path interpretation, homotopy would mean the existence of a continuous map
\[
p_{\bullet}\colon[0,1]\longrightarrow\bigl(M_n(\C[G]),\tau\bigr),\qquad p_0=p,\quad p_1=q,
\]
whose values are projections throughout. With these conventions fixed, a candidate formulation is: for every group \(G\) in the intended class, every integer \(n\geq1\), and all projections \(p,q\in M_n(\C[G])\), does such a homotopy imply that \(q=upu^{-1}\) for some \(u\in\operatorname{GL}_n(\C[G])\)?

This is a conditional formulation, not a uniquely determined reconstruction. The source does not specify the intended class of groups, the topology or algebraic meaning of homotopy, whether paths lie in the group ring or in a completion, whether stabilization is allowed, or the algebra in which the conjugating matrix must be invertible. These choices must be recovered from the source context before a definitive mathematical statement can be accepted.
\subsection{Short English statement}
The question asks whether projections connected by an appropriate homotopy must be conjugate by an invertible matrix over the complex group ring. The supplied text leaves the meaning of projection and homotopy, the group hypotheses, and possible stabilization unspecified. The displayed formulation therefore records a conditional reading requiring human mathematical review.
\subsection{Sources}
\begin{itemize}
\item Supplied local UnsolvedMath registry record: permanent ID \(30001332\), title ``Similarity of Homotopic Group-Ring Projections,'' source problem code \texttt{catalogue entry}, in the Oberwolfach Reports Open Problems collection. No individual proposer is supplied; the proposed year is recorded as \(2009\).
\item Supplied report citation: \emph{Noncommutative Geometry (2010)}, Oberwolfach Reports, DOI \texttt{10.4171/owr/2009/41}, pp.\ 2271--2334. Supplied source URL: \url{https://doi.org/10.4171/owr/2009/41}. The background describes Oberwolfach Reports as documenting mathematical workshops and their open-problem discussions.
\item Supplied literature source, labelled ``original Oberwolfach report 4084'': \url{https://doi.org/10.4171/owr/2009/41}.
\item Literature assessment checked at \texttt{2026-08-21}: the supplied status is ``open.'' The supplied assessment says that the extracted record contains only a generic conjecture marker and a statement that the topic is current research, and identifies no theorem resolving the underlying claim. This is supplied dated curation derived from the OWR material, not an independently verified literature assessment; time-sensitive status and cited claims require independent verification.
\item The registry labels the statement ``exact'' and reports that its mathematical content was checked against the official report, with notation, context, and question form normalized without changing the source-level claim. This verification claim is supplied provenance and has not been independently checked for this draft.
\end{itemize}
\end{document}
